% from ega0-1-fr.tex
\oldpage[0\textsubscript{I}]{77}

homomorphisme continu $w$ par passage au séparé complété. Or, si
$\mathfrak M$ est un idéal ouvert de $D$, il existe par hypothèse des idéaux
ouverts $\mathfrak K\subset B$, $\mathfrak L\subset C$ tels que
$u(\mathfrak K)\subset\mathfrak M$, $v(\mathfrak L)\subset\mathfrak M$~;
l'image par $w_0$ de
$\Im(\mathfrak K\otimes_A C)+\Im(B\otimes_A\mathfrak L)$ est encore
contenue dans $\mathfrak M$, d'où notre assertion.
\end{env}

\begin{proposition}[7.7.7]
\label{0.7.7.7-fr}
Si $B$ et $C$ sont deux $A$-algèbres préadmissibles,
$(B\otimes_A C)^\wedge$ est admissible, et si $\mathfrak K$ (resp.
$\mathfrak L$) est un idéal de définition de $B$ (resp. $C$), l'adhérence
dans $(B\otimes_A C)^\wedge$ de l'image canonique de
$\mathfrak H=\Im(\mathfrak K\otimes_A C)+\Im(B\otimes_A\mathfrak L)$ est
un idéal de définition.
\end{proposition}

Il suffit de montrer que $\mathfrak H^n$ tend vers $0$ pour la topologie
produit tensoriel, ce qui résulte immédiatement de l'inclusion
\[
  \mathfrak H^{2n}\subset
  \Im(\mathfrak K^n\otimes_A C)+\Im(B\otimes_A\mathfrak L^n).
\]

\begin{proposition}[7.7.8]
\label{0.7.7.8-fr}
Soient $A$ un anneau préadique, $\mathfrak J$ un idéal de définition de
$A$, $M$ un $A$-module \emph{de type fini}, muni de la topologie
$\mathfrak J$-préadique. Pour toute $A$-algèbre topologique adique et
noethérienne $B$, $B\otimes_A M$ s'identifie au produit tensoriel complété
$(B\otimes_A M)^\wedge$.
\end{proposition}

Si $\mathfrak K$ est un idéal de définition de $B$, il existe par hypothèse
un entier $m$ tel que $\mathfrak J^mB\subset\mathfrak K$, donc
$\Im(B\otimes_A\mathfrak J^{nm}M)=
\Im(\mathfrak J^{nm}B\otimes_A M)\subset
\Im(\mathfrak K^nB\otimes_A M)=\mathfrak K^n(B\otimes_A M)$~; on en
conclut que sur $B\otimes_A M$, le produit tensoriel des topologies de $B$
et de $M$ est la topologie $\mathfrak K$-préadique. Comme $B\otimes_A M$
est un $B$-module de type fini, la proposition résulte aussitôt de (7.3.6).

\subsection{Topologies sur les modules d'homomorphismes}
\label{subsection:0.7.8-fr}

\begin{env}[7.8.1]
\label{0.7.8.1-fr}
Soient $A$ un anneau $\mathfrak J$-adique noethérien, $M$ et $N$ deux
$A$-modules de type fini, munis de la topologie $\mathfrak J$-préadique~;
on sait (7.3.6) qu'ils sont séparés et complets~; en outre, tout
$A$-homomorphisme $M\to N$ est automatiquement continu, et le $A$-module
$\operatorname{Hom}_A(M,N)$ est de type fini. Pour tout entier $i\geq 0$, posons
$A_i=A/\mathfrak J^{i+1}$, $M_i=M/\mathfrak J^{i+1}M$,
$N_i=N/\mathfrak J^{i+1}N$~; pour $i\leq j$, tout homomorphisme
$u_j:M_j\to N_j$ applique $\mathfrak J^{i+1}M_j$ dans
$\mathfrak J^{i+1}N_j$, donc donne par passage aux quotients un
homomorphisme $u_i:M_i\to N_i$, ce qui définit un homomorphisme canonique
$\operatorname{Hom}_{A_j}(M_j,N_j)\to\operatorname{Hom}_{A_i}(M_i,N_i)$~; en outre, les
$\operatorname{Hom}_{A_i}(M_i,N_i)$ constituent un \emph{système projectif} pour ces
homomorphismes, et il résulte de (7.2.10) qu'il y a un isomorphisme
canonique
\[
  \varphi:\operatorname{Hom}_A(M,N)\longrightarrow
  \varprojlim_i\operatorname{Hom}_{A_i}(M_i,N_i).
\]
En outre~:
\end{env}

\begin{proposition}[7.8.2]
\label{0.7.8.2-fr}
Si $M$ et $N$ sont des modules de type fini sur un anneau $\mathfrak J$-adique
noethérien $A$, les sous-modules $\operatorname{Hom}_A(M,\mathfrak J^{i+1}N)$ forment un
système fondamental de voisinages de $0$ dans $\operatorname{Hom}_A(M,N)$ pour la
topologie $\mathfrak J$-adique, et l'homomorphisme canonique
\[
  \varphi:\operatorname{Hom}_A(M,N)\longrightarrow
  \varprojlim_i\operatorname{Hom}_{A_i}(M_i,N_i)
\]
est un isomorphisme topologique.
\end{proposition}

On peut en effet considérer $M$ comme quotient d'un $A$-module libre $L$ de
type fini, et par suite identifier $\operatorname{Hom}_A(M,N)$ à un sous-module de
$\operatorname{Hom}_A(L,N)$~; dans cette identification,
$\operatorname{Hom}_A(M,\mathfrak J^{i+1}N)$ est l'intersection de $\operatorname{Hom}_A(M,N)$ et de
$\operatorname{Hom}_A(L,\mathfrak J^{i+1}N)$~; comme la topologie induite sur
$\operatorname{Hom}_A(M,N)$ par la topologie $\mathfrak J$-adique de $\operatorname{Hom}_A(L,N)$ est
la topologie $\mathfrak J$-adique (7.3.2), on est ramené à démontrer la
première assertion

